\chapter{CTF 4 --- Vulnerability Assessment}
\label{ctf:ctf4}

\begin{notebox}[Target overview]
Vulnerability assessment CTF requiring identification and exploitation of web
application vulnerabilities across multiple attack surfaces. Each flag provided
a hint for the next.
\end{notebox}

% ─────────────────────────────────────────────────────────────────────────────
\section*{Tools Used}
\begin{itemize}
  \item \tool{nmap} --- port scanning and service detection
  \item \tool{dirb} / \tool{gobuster} --- directory enumeration
  \item \tool{curl} --- manual HTTP requests
  \item Browser --- direct inspection of discovered paths
\end{itemize}

% ─────────────────────────────────────────────────────────────────────────────
\section*{Flag 1 --- Exposed Version Control Repository}

\textbf{Hint:} \textit{Explore hidden directories for version control artifacts
that might reveal valuable information.}



\begin{termbox}
nmap -sC -sV target_ip
# Found /.git/ repository on port 80
# Browse to: http://target_ip/.git/
\end{termbox}

\begin{notebox}[Target overview]

The nmap \ic{-sC} default script scan found a \ic{.git} repository exposed
on port 80. Opened it in the browser and found flag 1.

\begin{takebox}
This is the second time a \ic{.git} exposure appeared in these CTFs. In real
engagements this is a critical finding --- the entire version history of the
application is exposed, including any secrets ever committed.
\end{takebox}

% ─────────────────────────────────────────────────────────────────────────────
\section*{Flag 2 --- Database via phpMyAdmin}

\textbf{Hint:} \textit{The data storage has some loose security measures. Can
you find the flag hidden within it?}

\end{notebox}

\begin{termbox}
# Check robots.txt first
curl http://target_ip/robots.txt
# Found: /phpmyadmin

# Browse to http://target_ip/phpmyadmin
# Login (default or weak credentials)
# Search database tables for flag
\end{termbox}

\begin{takebox}

\ic{robots.txt} revealed the phpMyAdmin path. Logged in, browsed the database
tables, and found the flag stored as a value in one of the tables.

% ─────────────────────────────────────────────────────────────────────────────
\section*{Flag 3 --- phpinfo.php Server Information}

\textbf{Hint:} \textit{A PHP file that displays server information might be worth
examining. What could be hidden in plain sight?}

\begin{termbox}
# Deeper script scan to find hidden files
nmap -sV --script=http-enum -p 80 192.232.28.3
# Found: /phpinfo.php

# Also discoverable via robots.txt (already checked above)
# Browse to: http://target_ip/phpinfo.php
# Search the page for the flag
\end{termbox}

The \ic{nmap --script=http-enum} scan found \ic{/phpinfo.php}. The phpinfo
page exposes server configuration, PHP settings, and environment variables.
The flag was embedded in the output.

\begin{notebox}
\ic{phpinfo.php} in production is a significant security risk. It reveals PHP
version, server software, file paths, environment variables (including any
hardcoded credentials), and loaded extensions. Flag or not, this is a finding
worth documenting in a real report.
\end{notebox}

% ─────────────────────────────────────────────────────────────────────────────
\section*{Flag 4 --- Sensitive Directory from robots.txt}

\textbf{Hint:} \textit{Sensitive directories might hold critical information.
Search through carefully for hidden gems.}

\begin{termbox}
# Already found this from robots.txt check in flag 2
# robots.txt revealed: /passwords

curl http://target_ip/passwords
# Or browse directly: http://target_ip/passwords
\end{termbox}

The same \ic{robots.txt} check that revealed phpMyAdmin also listed a
\ic{/passwords} directory. Opened it directly --- flag 4 was there.

% ─────────────────────────────────────────────────────────────────────────────
\section*{Flags Summary}

\begin{checkbox}
\begin{itemize}
  \item[$\checkmark$] \textbf{Flag 1} --- Exposed \ic{/.git/} (nmap \ic{-sC})
  \item[$\checkmark$] \textbf{Flag 2} --- phpMyAdmin database access
        (\ic{robots.txt} $\rightarrow$ login $\rightarrow$ table search)
  \item[$\checkmark$] \textbf{Flag 3} --- \ic{/phpinfo.php}
        (nmap \ic{--script=http-enum})
  \item[$\checkmark$] \textbf{Flag 4} --- \ic{/passwords} directory
        (\ic{robots.txt})
\end{itemize}
\end{checkbox}

% ─────────────────────────────────────────────────────────────────────────────
\section*{Overall Lessons from All CTFs}

\begin{takebox}
\textbf{Patterns that appeared across every CTF:}
\begin{itemize}
  \item \ic{robots.txt} gave flags or hints in three out of four CTFs.
        It is always the first check.
  \item \ic{nmap -sC} (default scripts) found hidden directories and
        exposed repositories that manual browsing would have missed.
  \item phpMyAdmin with weak credentials appeared twice. Default/weak credentials
        on admin interfaces is a recurring real-world vulnerability.
  \item CTF flags leave hints for the next flag. Read the flag content carefully
        before moving on.
  \item Exposed \ic{/.git/} repositories appeared twice. This is a common
        developer mistake that has real-world consequence.
\end{itemize}

\textbf{My CTF checklist (in order):}
\begin{enumerate}
  \item Port scan: \ic{nmap -sC -sV target}
  \item Check \ic{/robots.txt}
  \item Check \ic{/phpinfo.php}, \ic{/.git/}, \ic{/phpmyadmin}
  \item Deep web scan: \ic{nmap --script=http-enum -p 80 target}
  \item Directory brute-force: \ic{gobuster dir -u http://target -w wordlist}
  \item For each found path: browse, check for flags, note hints
\end{enumerate}
\end{takebox}

\end{takebox}